\documentclass[11pt,a4paper]{article}
\usepackage[utf8]{inputenc}
\usepackage[T1]{fontenc}
\usepackage[spanish]{babel}
\usepackage{hyperref}
\usepackage{geometry}
\geometry{margin=2.5cm}
\usepackage{parskip}

\title{344 las preguntas que obligan a ver como uno}
\author{Alejandro Valenzuela Rivera\\\small Universidad Autónoma de Nuevo León}
\date{\today}
\begin{document}
\maketitle

\begin{quote}
\textbf{Sesgo declarado.} Este trabajo se sostiene desde el sesgo declarado de su autor, y no finge neutralidad. Lo que afirma se puede tumbar, y ese es el punto: el falsificador está escrito, no escondido.
\end{quote}

\section*{Resumen}
Describir lo humano es un proceso infinito: ni nosotros sabemos qué significa. Y lo que se define desde un solo sitio obliga a ver como el que define.

\section{Por qué se sostiene}
Un constructo social definido desde un solo lugar es la trampa de AM: creer que todos necesitan lo mismo.

\section{De dónde viene}
La autoridad: 'lo humano se deja en blanco a propósito'.

\section{El caso}
Lo humano en blanco: el lugar que se deja sin definir para que el otro lo llene.

\section{Método}
\textquestiondown{}La definición deja al otro definir, o le impone una vista?

\section{Límites, y qué lo tumba}
Una definición de lo humano que no obligue a ver como el que la escribió.

\section{Lo que queda abierto}
\textquestiondown{}Se puede decir algo sobre lo humano sin obligarlo a ver como uno?

\section*{Cómo se leyó esta tesis}
Esto se escribió en el repositorio man-made-empathy, del que este candado es una pieza: candados/02 el alma y el libre albedrio/344 las preguntas que obligan a ver como uno.md. La máquina que lo compuso lo hizo desde este mismo archivo, y no agregó nada que no estuviera escrito. Si algo de esto es falso, se dice en el mismo sitio, y se corrige con fecha.

\end{document}
